% Values: ILLUMINATING
% Evidence planned:
% - ILLUMINATING: the case for a formal definition of Elm (implementers, a
%   second implementation in eco, programmers who reason about programs); the
%   character of the language as it shapes the definition (D2 to D7, D10,
%   D11); the method, with the keystone sentence B |- P => M stated and
%   explained in words, rules from constituent phrases, and the factoring
%   into static and dynamic semantics; the debt to the Definition of
%   Standard ML (D1) and the two conventions that the absence of state and
%   exceptions removes; what is left to implementations and where rulings
%   declare existing behaviour non-conforming (Appendix E).

\chapter*{Preface}\addcontentsline{toc}{chapter}{Preface}

Until now the meaning of an Elm program has been settled in practice by what
one implementation does with it: the Elm compiler for JavaScript, together
with the JavaScript kernel of its core library. The language is described in
prose, in a guide and in the documentation of its libraries, and that
description serves a programmer learning the language well. It serves less
well three other readers. An implementer must decide, for every construction,
what the program means, and cannot always tell from the existing
implementation whether a behaviour is the language's or the accident of a
host. A programmer who wants to show that two programs are equivalent, or
that a function never fails, needs a meaning precise enough to argue from.
And two parties who must agree on what a given program does need a text that
both accept as the authority.

This book is a formal definition of both the grammar and the meaning of Elm,
version 0.19.1. The occasion for it is a second implementation, eco, which
compiles Elm to native code and has its own runtime. Two implementations on
different hosts give different answers wherever an answer was a consequence
of the host rather than of the language: the range of integers, what a
character is, what happens when an integer is divided by zero. Once there
are two, each such question must be answered once, independently of both,
and precisely enough that either implementation can be checked against the
answer. We call the language so defined \emph{Standard Elm}.

\section*{Standard Elm}

Elm is a functional language for writing interactive programs that run
under the control of a host, first of all a web browser. It is pure: a
phrase denotes a value, and evaluating it changes nothing, since the
language has no assignable variables, no mutable data and no input or
output. It is strict: the argument of a function is evaluated before the
function is applied. It is statically typed, with types inferred and
polymorphism introduced at definitions. In its semantics it is extremely
close to ML: values, closures, custom types with constructors, pattern
matching and polymorphic type inference behave as they do there, and where
the two languages agree this definition gives Elm the meaning ML has. % D7

Five features distinguish Elm from ML and shape the definition. First, Elm
has no exceptions and pursues the aim that a program which has been accepted
does not fail when it runs. No phrase raises or handles an exception, and
the definition gives every phrase that has elaborated either a value or
none. A phrase has no value only when its evaluation does not terminate,
when it needs a top-level value while that value is itself being evaluated,
or when it reaches \kw{Debug.todo}, the one escape the language provides.
% D2, A25, A62
Arithmetic is accordingly total: where an operation has no mathematically
correct result, the definition gives a defined answer in its place, so that
\kw{3 // 0} is \kw{0}. % D3, A42
The basic arithmetic and comparison operations are defined in this book
rather than left to a library. % D4
Second, records are row-polymorphic: a function may demand some fields of a
record and accept any record that has them, and a record type may end in a
row variable standing for the fields not named. % D5
Third, a type variable may be constrained to a class (\kw{number},
\kw{comparable}, \kw{appendable} or \kw{compappend}) and then stands only for
the few predefined types on which arithmetic, ordering or appending is
defined. The classes are fixed: a program cannot add types to them. They
take the place of ML's overloaded identifiers and equality types, and
equality itself is defined at every type. % A35
Fourth, the module system is deliberately simple. A module is a sequence of
top-level declarations and makes some of them available to other modules by
name; there are no signatures, functors or sharing. % D6
Fifth, effects are managed. A program never performs an effect. It computes
values that describe effects (\emph{commands}, \emph{subscriptions} and
\emph{tasks}) and hands them to the platform that runs it; the platform
performs them and returns their outcomes to the program as messages, one
at a time. Commands, subscriptions and tasks are defined as types of the
basis rather than as syntax, but they are integral to the language and have
a chapter of their own. % D10
A program communicates with its host through \emph{ports}, which it
declares; \kw{port} is a reserved word, port declarations belong to the
language, and the definition states exactly which values can and cannot be
sent over a port. % D11

Standard Elm leaves out what Elm reserves to its own core packages: effect
modules, declarations of new infix operators and references to kernel code.
It also leaves out embedded shader blocks, which belong to a graphics
extension rather than to the language.

\section*{Language Definition}

The words of a programming language (\kw{type}, \kw{case}, \kw{port},
\kw{module}) name things that are not phrases: types, values, functions,
interfaces, commands. These things are the \emph{semantic objects} of the
language, and they, not its syntax, carry its character; when programmers
think about an Elm program they think in terms of them. For Standard Elm the
semantic objects include types with rows and classes, type schemes,
environments and module interfaces on the one hand, and values, closures,
records, tasks, commands and the boundary values (JSON values, with byte
sequences) that cross a port on the other. % A92
A definition of the language must therefore be a definition in terms of
these objects. They must be described in a notation independent of Elm,
since to describe them in another programming language would only move the
question to that language; ordinary mathematics serves. And they must be
given enough theory (substitution, instantiation, closure, the properties on
which the rules rely) for the rules to make sense, as a mathematician
develops the theory of the objects a subject is about.

With the objects in hand, the definition must say what meaning each phrase
of the language has. Here the keystone of our method is that a phrase never
has a meaning in isolation. It always has a meaning against a background,
which is itself a semantic object: a distillation of what earlier phrases
have contributed, such as the types and values of the names in scope and the
interfaces of the modules imported. The meaning of a phrase depends on its
background as much as on the phrase. The basic assertion of the definition
is therefore a sentence of the form
\[
  \jdg{B}{P}{M},
\]
read ``against the background $B$, the phrase $P$ has the meaning $M$''. The
formal purpose of this definition is exactly to decree which sentences of
this form hold, and no more. It does so by rules, each of which allows a
sentence about a compound phrase to be inferred from sentences about its
constituent phrases, so that the meaning of a phrase is determined by the
meanings of its parts. The rules are numbered through the whole book, and
each group of them is followed by Comments on the points that are not
obvious.

The meaning of an Elm phrase has two parts: what can be known of it before it
is evaluated, namely whether it is well formed and what its type is, and the
value it evaluates to. We give the two separately. The \emph{static
semantics} infers types and interfaces, and the \emph{dynamic semantics}
infers values; each is an instance of the keystone, with its own backgrounds
and its own meanings. The factoring is genuine. A value never has to be
presented together with its type, and the dynamic semantics consults no type
except where a port must know the type of what it sends or receives. A
phrase that does not elaborate has no meaning at all, and the dynamic
semantics need not consider it.

\section*{This Definition}

This book follows \emph{The Definition of Standard ML (Revised)} closely, in
its structure, its notation and its style: a grid of phases against levels,
a bare language with derived forms, semantic objects defined in figures,
numbered inference rules with Comments, and an initial basis in appendices.
% D1
Where Elm agrees with ML we have followed the treatment of that book, so
that a reader who knows it will find most of the static semantics of the
Core, and much of its dynamic semantics, in familiar form. % D7
The additions are the row variables and classes of the static semantics,
the much smaller static semantics of modules, a chapter on tasks, commands,
subscriptions and ports, and a chapter on programs which defines how the
platform runs a program as a labelled transition system whose observable
behaviour is its traffic on ports. Appendix~\ref{ch:appendix-standard-ml}
compares the two definitions construct by construct.

Two conventions of the Definition of Standard ML are absent here. There,
every sentence of the dynamic semantics carries a state, for references, and
the state is threaded through the premises of every rule, so that the order
of the premises is the order of evaluation; and for every rule there are
further implicit rules by which an exception packet raised in a premise
becomes the result of the conclusion. Standard Elm has neither state nor
exceptions, so neither convention is needed. Its dynamic sentences carry no
state, the order of premises carries no meaning, and no rule produces a
packet. % D2, A58

The body of the definition does not settle everything. Where implementations
legitimately differ (the range of integers, what a character is, the
treatment of floating-point special values, the order in which the commands
of a batch are performed), the body states only what every implementation
must satisfy, through abstract parameters, and
Appendix~\ref{ch:appendix-implementation} records the choices of Elm on
JavaScript and of eco, with their consequences for programs. % D8
Elsewhere the two implementations disagreed, or agreed on a behaviour that
followed from their hosts rather than from the language, and the definition
had to decide. Each such decision is a ruling, taken in the spirit of the
language: no run-time errors, total arithmetic, and a meaning that does not
depend on how a program is optimised. Some rulings declare existing
behaviour non-conforming. On Elm on JavaScript, \kw{modBy 0 x} stops the
program and comparing two functions with \kw{==} stops it too, where the
definition gives \kw{0} and a Boolean; on eco, a negative power of
\kw{-1} or \kw{1} is \kw{0} where the definition gives \kw{1} or \kw{-1};
on both, a value containing a NaN may be equal to itself, where the
definition says that it is not. % A42, A43, A50, A51
Appendix~\ref{ch:appendix-implementation} lists every such point with the
ruling that settles it.

\section*{How this Definition was Produced}

This definition was written with the assistance of an AI system, Claude
(model Claude Opus 5.5, made by Anthropic), working within a staged,
documented workflow. Under the author's direction it surveyed the sources
and test suites of the two implementations, Elm 0.19.1 and eco, recording
each fact with a citation to the code or test that establishes it; it
proposed a ruling, with its reasons, wherever the implementations disagree
or are silent; it drafted the chapters and appendices; it audited the
drafts, in independent passes, for consistency between chapters; and it
encoded the dynamic semantics of the Core in the proof assistant Lean~4,
namely every rule of Chapter~\ref{ch:dynamic-core} and the arithmetic,
equality, order, Boolean, conversion, elementary-function, append and
character part of $\APPLY$, with an evaluator proved sound
with respect to those rules and checked against the expected results of
487 of the test programs that accompany this book, which a program translates
into the encoding from their source. The author set the directions the
definition follows, reviewed every ruling, and is responsible for the
whole; the AI system is not an author. The survey, the record of rulings,
the audit reports and the mechanisation are kept with the sources of this
book.

\section*{Acknowledgements}

Elm was designed by Evan Czaplicki, and the language defined here is his.
The form of this book is that of \emph{The Definition of Standard ML} by
Robin Milner, Mads Tofte, Robert Harper and David MacQueen; we are indebted
to them for a method which has served a second language with little change.

\bigskip
\begin{flushright}
Rupert Smith\\[2ex]
November 2026
\end{flushright}
